% !PRINTED 702
\weberpage{702}
%<<CONT
associirten durch $\varphi_1$ theilbaren ganzen Zahlen, deren Normen nicht grösser als $t_1$ sind, und nach dem Satze \S.\,190,\,1. ist
\weq{2}{\mathop{\mathrm{Lim}}\limits_{t_1=\infty}\frac{T_1}{t_1} = \frac{g}{N_a(\varphi_1)},}
wenn $g$ die an der erwähnten Stelle angegebene Bedeutung hat, also eine \gesperrt{von Null verschiedene}, durch die Natur des Körpers $\Omega$ bestimmte Zahl ist. Wenn jetzt $T_2, t_2, \ldots, T_h, t_h$ die entsprechende Bedeutung für die Classen $A_2, \ldots, A_h$ haben, wie $T_1, t_1$ für $A_1$, so ist
\[
T = T_1 + T_2 + \cdots + T_h,
\]
\[
\frac{T}{t} = \frac{T_1}{t_1}\, N_a(\varphi_1) + \frac{T_2}{t_2}\, N_a(\varphi_2) + \cdots + \frac{T_h}{t_h}\, N_a(\varphi_h),
\]
und aus (1) und (2) ergiebt sich die fundamentale Formel:
\weq{3}{\mathop{\mathrm{Lim}}\limits_{s=1}\sum \frac{s-1}{N(\mathfrak{a})^s} = gh.}
% [sic] The following source citation is section 185,(11); the formula defining g is section 190,(11) on printed p. 693.
Die Zahl $g$ können wir nach \S.\,185,\,(11) als bekannt betrachten, wenn auch ihre wirkliche Berechnung noch auf der Voraussetzung beruht, dass ein Fundamentalsystem von Einheiten bekannt sei, und wenn auch die Ermittelung eines solchen Systemes in höheren Körpern immer als eine der grössten Schwierigkeiten betrachtet worden ist. Dann hängt die Berechnung der Classenzahl noch von der Bestimmung des Grenzwerthes auf der linken Seite von (3) ab, die natürlich auch nur in besonderen Fällen gelingt, doch aber oft wichtige Schlüsse über die Natur der Classenzahl gestattet. Wir wollen noch eine die Berechnung vorbereitende Umformung der Summe
\weq{4}{\sum \frac{1}{N(\mathfrak{a})^s} = \Phi(s)}
in ein unendliches Product entwickeln.

Es sei $\mathfrak{p}$ irgend ein Primideal $f^{\mathrm{ten}}$ Grades im Körper $\Omega$, also
\[
N(\mathfrak{p}) = p^f.
\]
Dann ist die unendliche geometrische Reihe
\[
1 + \frac{1}{N(\mathfrak{p})^s} + \frac{1}{N(\mathfrak{p})^{2s}} + \cdots = \frac{1}{1-p^{-sf}};
\]
wenn man diese Reihen für alle verschiedenen Primideale $\mathfrak{p}$ mit einander multiplicirt, so erhält man nach bekannten Sätzen aus
%CONT>>
